\scanpage{172}

gives

\deleted{\(\unsure{g_n\geq 3}\)}
\[
\underbrace{\underbrace{(2-2g_n)}_{\text{EP of }\widehat U_n}-3n}_{\text{EP}_{!}\text{ of }U_n}
=n^2\left(\underbrace{\underbrace{2}_{\text{EP of }\unsure{\widehat U_1}}-3}_{\text{EP of }U_1}\right)
\]

i.e.
\[
\inserted{2(g_n-1)=n^2-3n}
\qquad
\deleted{\unsure{2-2g_n=-n^2}}
\qquad \text{i.e.}
\]
\[
\text{(369)}\qquad g_n=\deleted{\unsure{2+}}1+\frac12 n(n-3)
\]

So we have
\[
\text{(370)}\qquad g_1=g_2=0,\quad g_3=1,\quad g_n>1\quad\text{if }n\geq\deleted{\unsure{3}}\inserted{4}.
\]

In any event, all the $U_n$ are anabelian, and even the $\widehat U_n$ for $n\geq4$.

If the \inserted{``\uline{anabelian}''} yoga developed in the first part of these notes is valid, the \deleted{\unsure{existence}} \deleted{(or the nonexistence) of} rational/$\mathbb Q$ points of a curve $U_n$ \deleted{\unsure{are explained by|are expressed by}} \inserted{correspond} to the \uline{splittings} of the extension of profinite groups

\noindent (371)
\diagram{\begin{tikzcd}[column sep=small]
1 \arrow[r] & \Pi_1(U_{n\overline{\mathbb Q}}) \arrow[r] &
\underset{\inserted{\substack{\unsure{\cap}\\\Pi_1(\unsure{U_1})\simeq\unsure{M_{0,3}^{!}}}}}{\Pi_1(U_n)} \arrow[r] & \Gamma_{\mathbb Q} \arrow[r] & 1
\end{tikzcd}}

which are distinct from the ``splittings at infinity'' (localized at the \inserted{3n} points at infinity of $U_n$\ldots), or again, the splittings such that \inserted{the \uline{centralizer} of} $(\Gamma_{\mathbb Q}^{+}=\operatorname{Ker}(\Gamma_{\mathbb Q}\xrightarrow{\chi}\widehat{\mathbb Z}^{*}))$ in $\Pi_1(U_n)$ is reduced to $\{1\}$. So -- if the ``theorem

\scanpage{173}
\mspage{734}

of Fermat'' is indeed valid (for $U_n$, say), it should be \unsure{reflected|reflected} in the nonexistence of such a splitting. Besides, if we \unsure{succeeded} in proving the nonexistence of splittings of (371) other than those ``coming from the points at infinity'', the theorem of Fermat would be proved for $U_n$, even without having established the ``anabelian yoga'' (a proof of which seems quite out of reach!). Now it is by no means excluded that a sufficiently thorough analysis of $M_{0,3}$ and $N_{1,1}$ in the present direction -- also involving finer and finer ``approximations'' of $\Pi_1(U_n)$ and of the extension structure (371), might eventually make it possible to prove the nonexistence of ``nontrivial'' splittings of (371), for all $n\geq3$. Each new refinement of the ``Lie approximations'' of $M_{0,3}$ should lead to a reexamination of the situation (371) and of the question of its splittings, with a view to a possible proof of the ``theorems of Fermat''.

\scanpage{174}

It is imperative to start from Lie approximations at least as fine as \unsure{those|that} of the present \S, which appears as the ``minima'' approximation in which one can express ``algebraically'' the tower of Fermat curves, \deleted{\unsure{where}} as we saw in 13\textsuperscript{o}), starting from $M_{0,3}\subset M_{0,3}^{\text{ét}}$.

Taking into account the operations of $\mathfrak S_3$ on $U_n$, the exact sequence (371) moreover embeds in a richer exact sequence

\noindent (372)
\diagram{\begin{tikzcd}[column sep=small]
1 \arrow[r] & \Pi_1(U_{n\overline{\mathbb Q}}) \arrow[r] &
\underset{\inserted{\substack{\cap\\\Pi_1(U_1,\mathfrak S_3)\simeq M_{0,3}}}}{\Pi_1(U_n,\mathfrak S_3)} \arrow[r] & \Gamma_{\mathbb Q}\times\mathfrak S_3 \arrow[r] & 1,
\end{tikzcd}}

and the calculations made in the present \S\ seem to suggest that this richer extension structure will play an important role.

From the point of view of the Lie approximation of $M_{0,3}$ developed in the present \S, it is appropriate to replace the extension structure (\unsure{372}) by a less fine quotient structure, by dividing by an invariant subgroup of $\Pi_1(U_n,\mathfrak S_3)$ contained in

\scanpage{175}
\mspage{735}

$\Pi_1(U_n,\overline{\mathbb Q})$, namely the intersection in $M_{0,3}$ of this group with $[\widehat\Pi_{0,3},\widehat\Pi_{0,3}]$. We find a commutative diagram

\noindent (373)
\diagram{\begin{tikzcd}[column sep=small,row sep=large]
1 \arrow[r] & \widehat M[n] \arrow[r] \arrow[d,"\unsure{1}"'] &
\begin{gathered}
\inserted{\overbrace{\Pi_1(U_n,\mathfrak S_3)^{\sim}}^{\unsure{\mathcal E_n}}}\\[-2pt]
\deleted{\unsure{\widetilde G_{\rho,\sigma}[n]}}
\end{gathered}
\arrow[r] \arrow[d,"\unsure{\text{epimorphism}}"] & \Gamma_{\mathbb Q}\times\mathfrak S_3 \arrow[r,no head] \arrow[d,"\chi\times\mathrm{id}"] & 1 \\
1 \arrow[r] & \widehat M[n] \arrow[r,no head] \arrow[d,hook] &
\widetilde G_{\rho,\sigma}[n] \arrow[r] \arrow[d,hook] & \widehat{\mathbb Z}^{*}\times\mathfrak S_3 \arrow[r,no head] \arrow[d,"\unsure{1}"] & 1 \\
1 \arrow[r] & \widehat M \arrow[r] & \widetilde G_{\rho,\sigma} \arrow[r] & \widehat{\mathbb Z}^{*}\times\mathfrak S_3 \arrow[r] & 1
\end{tikzcd}}

where the image of $\widehat M[n]\longrightarrow\widehat M$ is $n\widehat M$.

Here is \deleted{\illegible{1}} a way of constructing ``algebraically'' the diagrams (373) (for variable $n$) from the diagram

\noindent (374)
\diagram{\begin{tikzcd}[column sep=small,row sep=large]
1 \arrow[r] & \widehat M \arrow[r] \arrow[d,"\unsure{1}"'] &
\begin{gathered}
\overbrace{\Pi_1(U_1,\mathfrak S_3)^{\sim}}^{\mathcal E_1}\\[-2pt]
\Vert\\[-2pt]
(M_{0,3}/[\widehat\Pi_{0,3},\widehat\Pi_{0,3}])
\end{gathered}
\arrow[r] \arrow[d] & \Gamma_{\mathbb Q}\times\mathfrak S_3 \arrow[r,no head] \arrow[d,"\chi\times\mathrm{id}"] & 1 \\
1 \arrow[r,no head] & \widehat M \arrow[r] &
\begin{gathered}
\widetilde G_{\rho,\sigma}\\[-2pt]
\simeq(\widehat{\mathbb Z}^{*}\times\mathfrak S_3).\widehat M
\end{gathered}
\arrow[r] & \widehat{\mathbb Z}^{*}\times\mathfrak S_3 \arrow[r] & 1.
\end{tikzcd}}

We choose a \deleted{\unsure{which}} splitting of the first i.e. a splitting of the second (two such splittings are conjugate by a unique element of $\widehat M$\ldots).
\margin{We have seen that these splittings \inserted{\unsure{correspond|correspond}} \deleted{\unsure{being}} bijectively to those of the extension $\mathfrak S^{+\wedge}/\mu$ of $\mathfrak S_3$ by $\widehat M$, and we have defined a canonical one \deleted{\unsure{in}} $12^\circ)$.}
\deleted{\unsure{This is moreover}} \inserted{It can be inter}preted as a \inserted{geometric} point of $\unsure{U_\infty|U_n}$ above the

\scanpage{176}

base point $-\overline j$ of $U_{0,3\mathbb Q}\simeq U_{1\mathbb Q}$ -- it defines a geometric point of each of the $U_{n\mathbb Q}$ above $-\overline j$. In terms of this splitting of $\mathcal E_1$, allowing us to write
\[
\text{(375)}\qquad \mathcal E_1\simeq(\Gamma_{\mathbb Q}\times\mathfrak S_3).\widehat M
\qquad\text{($\frac12$ direct product)},
\]
the $\underset{\inserted{\substack{\unsure{\rotatebox{90}{$\simeq$}}\\\Pi_1(U_n,\unsure{\mathfrak S_n|\mathfrak S_3})/\ldots}}}{\mathcal E_n}\subset\mathcal E_1$ are identified with the semidirect products
\[
\text{(376)}\qquad \mathcal E_n=(\Gamma_{\mathbb Q}\times\mathfrak S_3).
\underbrace{(n\widehat M)}_{\overset{\text{definition}}{=}\widehat M[n]}
\]

(the interpretations are at the level of generalities of the proposition-lemma on page~$712'$ ff, and would deserve to appear in the \unsure{scholium|scholium}-proposition in question\ldots)

A \unsure{points|point} \inserted{$x_n$} \unsure{rational|rational}/$\mathbb Q$ of $U_n$ defines a splitting of the extension (of $\Gamma_{\mathbb Q}$ by $n\widehat M=\widehat M[n]$)
\[
\text{(377)}\qquad \underset{\inserted{\unsure{\varphi}}}{\mathcal E_n^{!}}\simeq\Gamma_{\mathbb Q}.(n\widehat M)
\qquad\text{($\frac12$ direct)}
\]
quotient of (371), knowledge of which
\scanpage{177}
\mspage{736}

modulo \deleted{\unsure{$n$}} $\widehat M[n]$-conjugation \uline{implies} that
of the corresponding splitting of $\mathcal E_1$, hence
implies knowledge of the corresponding point \inserted{$x_1$}
of $U_1(\mathbb Q)$ \deleted{\unsure{$x_1$}}, which means
that $x_n$ is known modulo the operation of
$\unsure{\gamma_n}=\mu_n(\mathbb Q)^{\{0,1,\infty\}}/\mu_n(\mathbb Q)$.
Now we have
\[
\begin{gathered}
\text{(378)}\qquad
\unsure{\gamma_n}\overset{\mathrm{dfn}}=
\mu_n(\mathbb Q)^{\{0,1,\infty\}}/\mu_n(\mathbb Q)=1\\
\text{if }n\text{ odd}
\end{gathered}
\]
for $\mu_n(\mathbb Q)=1$ if $n$ odd -- so in
this case $x_n$ is known when one knows
the corresponding splitting, and even
\unsure{only} its $\widehat M[n]$-conjugacy class.

This leads us to determine all these
possible $\widehat M[n]$-conjugacy classes of
such splittings, in order to examine (if it can be
done) which ones can reasonably
correspond to an $x_n\in U_n(\mathbb Q)$. These classes
correspond to the elements of
\[
\begin{aligned}
\text{(379)}\qquad
H^1(\Gamma_{\mathbb Q},n\widehat M)
&\simeq H^1(\Gamma_{\mathbb Q},\widehat M)
\simeq
\underbrace{H^1\unsure{(}\Gamma_{\mathbb Q},\widehat{\mathbb Z}(-1))}
_{\unsure{(\mathbb Q^*)^\wedge|(\mathbb Q^*]^\wedge}}
\otimes_{\mathbb Z}M\\
&\simeq G(\mathbb Q)^\wedge
\end{aligned}
\]

\scanpage{178}

where
\[
G\overset{\mathrm{dfn}}=
\unsure{\mathbb G_{m\,\mathbb Q}|G_{m\,\mathbb Q}}^{\{0,1,\infty\}}
/\unsure{\mathbb G_m|G_m}
\]
is the torus of dimension $2$ over $\mathbb Q$ in which
the $U_n$ are embedded. We see that
everything happens as if we were studying
simply the existence (!) and the cohomological
locating of the $\mathbb Q$-rational points of
$G$ itself -- without concerning ourselves
with the equation $x_0^n+x_1^n+x_\infty^n=1$ which defines
$U_n$. This shows plainly that
the approximation on which we are working
is much too coarse -- in fact,
in practice we \unsure{remain} in the
domain of a typically
``abelian'' analysis. It is clear that it is
necessary to pass to finer
approximations.

\scanpage{179}
\mspage{737}

\uline{52}\qquad
\uline{Schematic envelopes \unsure{of|of the} L-groups.}

$1^\circ)$ \uline{Notion of L-group}, \uline{of LT-group} (discrete case).

The notion of L-group is the natural
generalization of that of ``group with loops'',
which is called for if we want to be able to encompass
mixed fundamental groups of
orientable surfaces (of ``finite \unsure{character}'', i.e.
homeomorphic to an $X\setminus S$, $X$ a compact surface,
$S$ a finite subset) with a finite group of automorphisms.

We have seen that if $\Pi$ is a
group with loops (discrete), then every subgroup
of finite index $\Pi_0$ inherits a
loop structure, taking as
``loop subgroups'' of $\Pi_0$ the intersections
with $\Pi_0$ of the loop \unsure{subgroups} of $\Pi$. [N.B. we
assume that we are not in the cases
$(g,\nu)$ \deleted{\unsure{with}} with $g=0$, $\nu\in\{0,1\}$, where the loop structure
is \inserted{not} determined by the sole datum of a
family of conjugacy classes of \unsure{subgroups}
\deleted{\unsure{discrete}} $(\simeq\mathbb Z)$ -- the \inserted{only} case that would cause a problem here

\scanpage{180}

is the case $(0,1)$ \inserted{of $\Pi\simeq\mathbb Z$} -- in the case $(0,0)$, we have $\Pi=\{1\}$
and $\Pi_0=\Pi\ldots$ ]. Conversely, the loop
structure of $\Pi$ is recovered from
that of $\unsure{\Pi_{0*}|\Pi_0}$ (when it exists!), the
loop \unsure{subgroups} of $\Pi$ being the centralizers
in $\Pi$ of the \unsure{loop-subgroups} of $\Pi_0$.
\deleted{The exceptional case}

Let $\Pi$ be a discrete group,
\inserted{\deleted{finitely generated}}
\deleted{let} and
let $\mathcal L_0(\Pi)$ be the set of group-with-loops
structures on $\Pi$, of type $(g,\nu)$
\deleted{\unsure{with|$\ne$}}\inserted{\unsure{anabelian} i.e. $\ne$}
$(0,\unsure{0|1}),
\deleted{\unsure{(0,2)}}\inserted{(0,1)},
\deleted{\unsure{(1,0)}}\inserted{(1,0)}$.
For the subgroups \deleted{of} $\Pi_i$ \inserted{of $\Pi$} of finite index,
the $\mathcal L_0(\Pi_i)$ form an inductive system
\deleted{\unsure{with a map|with maps}}
of sets. \deleted{We set} If $\Pi_i$ is normal,
\[
\deleted{\mathcal L(\Pi)\qquad
\unsure{\mathcal L(\Pi_i)|\mathcal L_0(\Pi_i)}}
\]
we \deleted{\unsure{can}} \unsure{set|take}
\[
\deleted{\text{(1)}\qquad
\mathcal L(\Pi)=\varinjlim_i\mathcal L_0(\Pi_i)}
\]
\[
\begin{gathered}
\text{(1)}\qquad
\mathcal L_0^{\unsure{\natural|\sharp}}(\Pi_i)=
\text{set of loop structures}\\
\text{on }\Pi_i\text{ invariant under}\\
G_i=\Pi/\Pi_i.
\end{gathered}
\]

\deleted{We assume \unsure{that} $\Pi$ admits a finite
generating set, \unsure{hence a system}}
If there exists a $\Pi_{i_0}$ of finite index such that
$\mathcal L(\Pi_{\unsure{i|i_0}})\ne\varnothing$, then $\Pi_{i_0}$ hence also $\Pi$ admits
a finite generating set, and it follows
that the set of \inserted{\uline{normal}} subgroups of finite index
is cofinal in the set of subgroups
of finite index. We shall then set